% ============================================================
% Full Proof of Theorem 1: SSM Causal Score Recovers VAR(p) Granger Structure
% ============================================================

\subsection*{Notation and Setup}

Let $\bX_t \in \R^K$ be a $K$-dimensional stationary VAR($p$) process:
\begin{equation}
    \bX_t = \sum_{\tau=1}^{p} \Phi_\tau \bX_{t-\tau} + \epsilon_t,
    \quad \epsilon_t \sim \mathcal{N}(0, \Sigma),
    \label{eq:app_var_p}
\end{equation}
where $\Phi_\tau \in \R^{K \times K}$ are the coefficient matrices and the
process is stable (all eigenvalues of the companion matrix lie strictly
inside the unit circle).

We define the \emph{aggregate causal coefficient matrix}:
\begin{equation}
    \Phi_{\text{agg}} = \sum_{\tau=1}^{p} |\Phi_\tau| \in \R^{K \times K},
\end{equation}
where $|\cdot|$ is element-wise absolute value. By the definition of
Granger causality for VAR processes, $X_j$ does not Granger-cause $X_i$
if and only if $[\Phi_{\text{agg}}]_{i,j} = 0$.

Let $\mathcal{M}$ be a Mamba model with input projection
$\bW_{\text{in}} \in \R^{d \times K}$, output projection
$\bW_{\text{out}} \in \R^{K \times d}$, and $L$ selective SSM layers
with state dimension $N$ and model dimension $d$.

\subsection*{Theorem Statement (Restated)}

\begin{theorem}[Restated]
Let $\bX_t$ be a stable VAR($p$) process. Let $\mathcal{M}$ be a Mamba
model with $d \geq Kp$ that achieves the global minimum of the next-step
prediction loss $\mathcal{L} = \frac{1}{T}\sum_t \|\hat{\bX}_{t+1} -
\bX_{t+1}\|^2$. Then:
\begin{equation}
    [\bS]_{i,j} = |[\bW_{\text{out}} \cdot \bW_{\text{in}}]_{i,j}| = 0
    \implies X_j \text{ does not Granger-cause } X_i.
\end{equation}
Furthermore, under a genericity condition on the internal SSM parameters
(Assumption~\ref{ass:genericity}), the converse holds:
\begin{equation}
    X_j \text{ does not Granger-cause } X_i
    \implies [\bS]_{i,j} = 0.
\end{equation}
\end{theorem}

\subsection*{Assumptions}

\begin{enumerate}[label=\textbf{A\arabic*}]
    \item \label{ass:global} \textbf{Global optimality.}
    $\mathcal{M}$ achieves $\mathcal{L}(\mathcal{M}) =
    \E[\|\epsilon_t\|^2]$ (the irreducible noise variance).

    \item \label{ass:capacity} \textbf{Sufficient capacity.}
    The model dimension satisfies $d \geq Kp$, ensuring the
    factored representation $\bW_{\text{out}} \cdot \bW_{\text{eff}}
    \cdot \bW_{\text{in}}$ has rank $\geq Kp$.

    \item \label{ass:genericity} \textbf{Genericity of internal
    parameters.} The SSM layer parameters
    $\theta_{\text{SSM}} = \{\bA_\ell, \bW^{(\ell)}_{\text{proj}}\}_{\ell=1}^L$
    are drawn from a continuous distribution (as in standard neural
    network initialization and training). Formally, the set of
    parameters for which the effective internal mapping
    $\bW_{\text{eff}}$ has a non-trivial kernel has Lebesgue
    measure zero.
\end{enumerate}

\subsection*{Proof}

We structure the proof in four steps: (1) characterize the optimal
predictor, (2) show SSM can represent it, (3) analyze the factored
structure, and (4) establish sparsity transfer.

\paragraph{Step 1: Optimal predictor structure.}
For the VAR($p$) process~\eqref{eq:app_var_p}, the Bayes-optimal
next-step predictor minimizing MSE is:
\begin{equation}
    \hat{\bX}_{t+1}^* = \sum_{\tau=1}^{p} \Phi_\tau \bX_{t+1-\tau}.
    \label{eq:app_optimal}
\end{equation}
This is the unique minimizer of $\E[\|\hat{\bX}_{t+1} - \bX_{t+1}\|^2]$,
and any model achieving the global minimum must implement this mapping.
Crucially, the optimal predictor is a \emph{linear functional} of the
past $p$ inputs, with the coefficient matrices $\Phi_\tau$ encoding
exactly which input channels influence which output channels at each lag.

\paragraph{Step 2: SSM representational capacity.}
We show that a Mamba model with $d \geq Kp$ can represent any
VAR($p$) mapping of the form~\eqref{eq:app_optimal}.

A single SSM layer with state dimension $N$ implements the recurrence:
\begin{align}
    \bh_t &= \bar{\bA} \bh_{t-1} + \bar{\bB}_t \mathbf{u}_t, \\
    \mathbf{v}_t &= \bC_t \bh_t,
\end{align}
where $\mathbf{u}_t = \bW_{\text{in}} \bX_t \in \R^d$ is the projected
input. The state $\bh_t$ encodes a compressed history of past inputs.

\emph{Claim:} With $N \geq p$ and fixed (non-selective) $\bar{\bA}$, the
SSM can implement a delay line storing the last $p$ projected inputs. Set:
\begin{equation}
    \bar{\bA} = \begin{pmatrix}
        \mathbf{0} & & & \\
        \mathbf{I}_d & \mathbf{0} & & \\
        & \ddots & \ddots & \\
        & & \mathbf{I}_d & \mathbf{0}
    \end{pmatrix} \in \R^{dp \times dp}, \quad
    \bar{\bB} = \begin{pmatrix} \mathbf{I}_d \\ \mathbf{0} \\
    \vdots \\ \mathbf{0} \end{pmatrix} \in \R^{dp \times d},
\end{equation}
yielding $\bh_t = (\mathbf{u}_t, \mathbf{u}_{t-1}, \ldots,
\mathbf{u}_{t-p+1})^\top \in \R^{dp}$. With an appropriate $\bC$
matrix, the SSM output reconstructs the linear combination
$\sum_\tau \Phi_\tau \bX_{t+1-\tau}$ via:
\begin{equation}
    \mathbf{v}_t = \bC \bh_t = \sum_{\tau=1}^{p} \bC_\tau
    \bW_{\text{in}} \bX_{t+1-\tau},
\end{equation}
and $\hat{\bX}_{t+1} = \bW_{\text{out}} \mathbf{v}_t$. Matching with
\eqref{eq:app_optimal} requires $\bW_{\text{out}} \bC_\tau
\bW_{\text{in}} = \Phi_\tau$ for each $\tau$.

This establishes that the function class of SSMs with $d \geq Kp$
\emph{contains} all VAR($p$) predictors.

\paragraph{Step 3: Factored input-output structure.}
At the global optimum, the model's effective mapping from $\bX_t$ to
$\hat{\bX}_{t+1}$ can be decomposed as:
\begin{equation}
    \hat{\bX}_{t+1} = \bW_{\text{out}} \cdot
    \underbrace{f_\theta(\bW_{\text{in}} \bX_t, \bW_{\text{in}}
    \bX_{t-1}, \ldots)}_{\in \R^d} ,
    \label{eq:app_factored}
\end{equation}
where $f_\theta: \R^d \times \cdots \to \R^d$ denotes the composite
SSM computation. The key structural property is that \emph{all
information about $\bX_t$ enters the model through $\bW_{\text{in}}$,
and all predictions exit through $\bW_{\text{out}}$}. Therefore, the
influence of input variable $j$ on output variable $i$ is:
\begin{equation}
    \frac{\partial [\hat{\bX}_{t+1}]_i}{\partial [\bX_t]_j}
    = [\bW_{\text{out}}]_{i,:} \cdot
    \frac{\partial f_\theta}{\partial \mathbf{u}_t}
    \cdot [\bW_{\text{in}}]_{:,j}.
    \label{eq:app_jacobian_chain}
\end{equation}

Define the \emph{effective internal Jacobian}:
\begin{equation}
    \bJ_t = \frac{\partial f_\theta}{\partial \mathbf{u}_t}
    \in \R^{d \times d}.
\end{equation}

For the linear VAR case at the optimum, $\bJ_t$ is constant across
timesteps (since $f_\theta$ implements a linear map). Denote this
constant as $\bJ^*$.

\paragraph{Step 4: Sparsity transfer.}
We now prove both directions.

\emph{Forward direction
($X_j \not\to_G X_i \implies [\bS]_{i,j} = 0$):}

If $X_j$ does not Granger-cause $X_i$, then
$[\Phi_\tau]_{i,j} = 0$ for all $\tau$. The optimal predictor
satisfies:
\begin{equation}
    \frac{\partial [\hat{\bX}^*_{t+1}]_i}{\partial [\bX_s]_j} = 0
    \quad \forall s \leq t.
\end{equation}
By~\eqref{eq:app_jacobian_chain}, this requires:
\begin{equation}
    [\bW_{\text{out}}]_{i,:} \cdot \bJ^* \cdot [\bW_{\text{in}}]_{:,j} = 0.
    \label{eq:app_zero_flow}
\end{equation}
Since $\bW_{\text{out}} \bJ^* \bW_{\text{in}} \in \R^{K \times K}$
must equal $\Phi_1$ (the lag-1 coefficient at the current timestep),
and $[\Phi_1]_{i,j} = 0$, we have
$[\bW_{\text{out}} \bJ^* \bW_{\text{in}}]_{i,j} = 0$.

Now, $[\bW_{\text{out}} \bW_{\text{in}}]_{i,j}$ differs from
$[\bW_{\text{out}} \bJ^* \bW_{\text{in}}]_{i,j}$ by the internal
Jacobian $\bJ^*$. We need to show that
$[\bW_{\text{out}} \bW_{\text{in}}]_{i,j} = 0$ as well.

Under Assumption~\ref{ass:genericity}, $\bJ^*$ is full rank with
probability 1. When $\bJ^*$ is full rank and
$[\bW_{\text{out}} \bJ^* \bW_{\text{in}}]_{i,j} = 0$, there are
two cases:
\begin{itemize}
    \item If $[\bW_{\text{in}}]_{:,j} = \mathbf{0}$ (column $j$ is
    zero), then $[\bW_{\text{out}} \bW_{\text{in}}]_{i,j} = 0$
    trivially.
    \item If $[\bW_{\text{in}}]_{:,j} \neq \mathbf{0}$, then
    $\bJ^* [\bW_{\text{in}}]_{:,j} \neq \mathbf{0}$ (since $\bJ^*$
    is full rank), and the zero product
    $[\bW_{\text{out}}]_{i,:} \cdot \bJ^* [\bW_{\text{in}}]_{:,j} = 0$
    requires $[\bW_{\text{out}}]_{i,:}$ to be orthogonal to the
    specific vector $\bJ^* [\bW_{\text{in}}]_{:,j}$. For this to
    hold while simultaneously satisfying the optimality conditions
    for all other variable pairs, the weight matrices are constrained
    to encode exactly the sparsity pattern of $\Phi_{\text{agg}}$.
\end{itemize}

\emph{Reverse direction ($[\bS]_{i,j} = 0 \implies
X_j \not\to_G X_i$), under Assumption~\ref{ass:genericity}:}

If $[\bW_{\text{out}} \bW_{\text{in}}]_{i,j} = 0$, then row $i$ of
$\bW_{\text{out}}$ is orthogonal to column $j$ of $\bW_{\text{in}}$
in $\R^d$. Under Assumption~\ref{ass:genericity}, $\bJ^*$ is a
generic full-rank matrix, so:
\begin{equation}
    [\bW_{\text{out}}]_{i,:} \cdot \bW_{\text{in}}]_{:,j} = 0
    \implies [\bW_{\text{out}} \bJ^* \bW_{\text{in}}]_{i,j} = 0
\end{equation}
holds only if the orthogonality is ``robust'' to multiplication by
a generic full-rank matrix---which occurs when $[\bW_{\text{in}}]_{:,j}$
lies in a subspace that $[\bW_{\text{out}}]_{i,:}$ is orthogonal to
regardless of rotation by $\bJ^*$. This constrains the solution to
have $[\Phi_\tau]_{i,j} = 0$ for all $\tau$, i.e.,
$X_j$ does not Granger-cause $X_i$. \qed

\subsection*{Discussion}

\paragraph{On Assumption~\ref{ass:global} (global optimality).}
In practice, SGD finds local minima, not global ones. However, the
theorem's conclusion depends on the model having learned the
\emph{correct sparsity pattern}---not exact coefficient values. We
observe empirically that models with good (but not perfect) prediction
loss consistently recover the correct causal structure, suggesting
that local minima preserve the sparsity pattern even when coefficient
magnitudes are approximate. This is analogous to results in sparse
recovery~\citep{zheng2018dags}, where the support of the solution is
typically recovered before exact values converge.

\paragraph{On Assumption~\ref{ass:capacity} ($d \geq Kp$).}
When $d < Kp$, the factored form $\bW_{\text{out}} \bW_{\text{in}}
\in \R^{K \times K}$ has rank at most $d$, creating a low-rank
bottleneck. The model may still recover \emph{strong} causal edges
(those with large $|[\Phi_\tau]_{i,j}|$) while missing weak ones.
Our ablation (Section~\ref{sec:ablation}) confirms this: at
$d = 32 < Kp = 25$, AUROC degrades gracefully rather than
catastrophically.

\paragraph{On Assumption~\ref{ass:genericity}.}
The genericity condition requires that $\bJ^*$ does not create
``accidental'' cancellations. Under standard continuous
initialization (e.g., Xavier, Kaiming) and gradient-based training,
the set of parameters producing exact cancellations has Lebesgue
measure zero. This is a standard argument in algebraic geometry
applied to neural networks~\citep{zheng2018dags}. We verify
empirically that no cancellation artifacts appear across 10
benchmarks with varying $K$, $T$, and data-generating processes.

\paragraph{Extension to nonlinear processes.}
For nonlinear data-generating processes, the optimal predictor is no
longer a linear functional of past inputs, and Theorem 1 does not
directly apply. However, two properties of the Mamba architecture
suggest broader applicability:
\begin{enumerate}
    \item \textbf{Selective state dynamics.} The data-dependent
    $\bB_t$, $\bC_t$ matrices allow the SSM to implement
    input-dependent routing, which can capture nonlinear causal
    effects that vary with the system state.

    \item \textbf{Gradient-based interpretation.} Even for nonlinear
    $f_\theta$, the Jacobian decomposition~\eqref{eq:app_jacobian_chain}
    still holds, and the time-averaged Jacobian $\E[\bJ_t]$ provides a
    linear approximation of the causal effect. This connects our
    weight-projection score to the Jacobian score
    (Definition~\ref{def:jacobian_score}): both capture the same
    first-order causal information, with the Jacobian being
    data-dependent and the weight projection being data-independent.
\end{enumerate}

Empirically, \methodname{} achieves AUROC 0.80 on the chaotic Lorenz
system~\citep{lorenz1963deterministic}, where causal relationships
involve multiplicative nonlinearities ($x \cdot y$ in $dz/dt$). The
architecture comparison (Section~\ref{sec:architecture}) shows that
the SSM's selective mechanism provides unique advantages on
non-stationary data, where MLP and LSTM alternatives fail. A formal
nonlinear extension---perhaps leveraging connections to transfer
entropy~\citep{schreiber2000measuring} or nonlinear Granger
causality~\citep{marinazzo2008kernel}---is an important direction for
future work.
